\documentclass[11pt]{article}
\usepackage[margin=0.95in]{geometry}
\usepackage[T1]{fontenc}
\usepackage{lmodern}
\usepackage{amsmath,amssymb,amsthm}
\usepackage[hidelinks]{hyperref}
\newtheorem*{claim}{Claim}
\newcommand{\C}{\mathbb C}

\begin{document}
\begin{center}
{\Large Literature-implied negative answer for arXiv:1407.1065}\\[4pt]
{\large The untrimmed Wirtinger Flow initializer cannot use $m=O(n)$ samples}
\end{center}

\paragraph{Status.}
\textbf{literature\_implied\_answer (full negative answer to the stated
linear-sampling question).}  This is not a new run counterexample.  A later
paper identifies the extreme-measurement obstruction; the short argument
below makes its implication for the source's $1/8$ basin explicit.

\paragraph{Source question.}
In Cand\`es, Li, and Soltanolkotabi, \emph{Phase Retrieval via Wirtinger Flow:
Theory and Algorithms}, arXiv:1407.1065, PDF page 6, Theorem 3.3 proves the
Gaussian-model guarantee with $m\geq c_0n\log n$.  Immediately afterward the
authors observe that injectivity needs only linearly many samples and ask
whether Theorem 3.3 holds when $m$ is merely proportional to $n$.

The theorem uses the leading eigenvector of
\[
 Y=\frac1m\sum_{r=1}^m y_ra_ra_r^*,\qquad
 y_r=|\langle a_r,x\rangle|^2,                                      \tag{1}
\]
scaled to have norm asymptotic to $\|x\|$, and requires
$\operatorname{dist}(z_0,x)\leq\|x\|/8$.

\paragraph{Supporting literature.}
Chi, Lu, and Chen, \emph{Nonconvex Optimization Meets Low-Rank Matrix
Factorization: An Overview}, arXiv:1809.09573, Section 8.3.1 (PDF page 40),
cites the standard phase-retrieval spectral construction and proves
\[
 \|Y\|\geq \frac{(\max_r y_r)\|a_{r_*}\|^2}{m}
       \asymp \frac{n\log m}{m}.                                    \tag{2}
\]
They conclude that this is incompatible with linear sample complexity and
that the standard recipe must be modified by truncating the large
measurements.  Their Theorem 27 records the resulting $m=O(n)$ guarantee for
the \emph{truncated} spectral estimator, not the source initializer.

\paragraph{Identification and rigorous bridge.}
Normalize $\|x\|=1$ and take $m/n\to\alpha\in(0,\infty)$.  Formula (2) tends
to infinity in probability (for complex Gaussian vectors, $y_r$ is
exponential; the same maximum argument applies).  On the other hand,
\[
 x^*Yx=\frac1m\sum_{r=1}^m y_r^2\longrightarrow \mathbb E y_1^2=2,    \tag{3}
\]
and either normalization in the source makes $\|z_0\|^2\to1$.

Let $\Lambda=\|Y\|$ and let $z_0$ be a scaled leading eigenvector.  If, after
a phase rotation, $z_0=x+h$ with $\|h\|\leq1/8$, positivity of $Y$ and
$|u+v|^2\leq2|u|^2+2|v|^2$ give
\[
 \|z_0\|^2\Lambda=z_0^*Yz_0
 \leq2x^*Yx+2h^*Yh
 \leq2x^*Yx+\frac1{32}\Lambda.                                    \tag{4}
\]
With probability tending to one, $\|z_0\|^2\geq1/2$ and $x^*Yx\leq3$,
so (4) would force $\Lambda\leq 64/5$.  This contradicts (2).  Thus
\[
 \Pr\{\operatorname{dist}(z_0,x)\leq1/8\}\longrightarrow0.          \tag{5}
\]

\begin{claim}
For every fixed sampling ratio $\alpha>0$, the untrimmed spectral initializer
in arXiv:1407.1065 fails the initialization conclusion of Theorem 3.3 with
probability tending to one when $m/n\to\alpha$.  Consequently Theorem 3.3
cannot hold in the proposed linear-sampling regime with that initializer.
\end{claim}

For completeness, (2) itself follows by choosing an index with
$y_r\geq\tfrac12\log n$ (such an index exists with probability tending to
one) and using $\|a_r\|^2\geq n/2$.  The probability of the first event is at
least $1-\exp(-\Theta(\sqrt n))$ and the second fails with exponentially small
probability.  Hence $\Lambda\geq c_\alpha\log n$.

\paragraph{Scope.}
This does not rule out linear-sample phase retrieval or modified Wirtinger
Flow.  Truncated Wirtinger Flow and other preprocessed spectral methods do
achieve $m=O(n)$.  It rules out exactly the source theorem in the proposed
regime with its original untrimmed spectral initialization.

\paragraph{Search evidence.}
The run indexes contained no prior packet for arXiv:1407.1065.  Exact-question,
initializer, truncation, and extreme-measurement searches located
arXiv:1809.09573.  Later phase-transition papers analyze bounded preprocessing
and likewise motivate truncation; no evidence was found that the raw
initializer was subsequently rescued at fixed aspect ratio.

\begin{thebibliography}{9}
\bibitem{CLS}
E. J. Cand\`es, X. Li, and M. Soltanolkotabi,
\emph{Phase Retrieval via Wirtinger Flow: Theory and Algorithms},
IEEE Trans. Inform. Theory 61 (2015), 1985--2007; arXiv:1407.1065.
\bibitem{CLC}
Y. Chi, Y. M. Lu, and Y. Chen,
\emph{Nonconvex Optimization Meets Low-Rank Matrix Factorization: An Overview},
IEEE Trans. Signal Process. 67 (2019), 5239--5269; arXiv:1809.09573.
\end{thebibliography}
\end{document}

